% =====================================================================
% Section 6 -- Application to field conveyors. Written in step 6.6.
% Inputs: notes 4.26-4.27; fig_applications; table_applications (generated
% by maps/paper_tables.py). Lodewijks (1996) re-read in step 6.6: drive
% grip checked with the start-up factor 1.2 of a drain-type fluid coupling
% (Table 8.2, Eqs. 8.7-8.9, pp. 154-155); only single hinge elements, so no
% slip, at the drive pulleys (p. 137). DIN 22101 sec. 8.2.3 re-read: relative
% sag below 0.01 in steady operation, more allowed in non-steady operation.
% Numbers: paper_numbers.py, section "Application to field conveyors".
% Tests: test_applications.py, test_tables.py.
% =====================================================================
\section{Application to field conveyors}\label{sec:application}

\begin{figure}[tbp]
  \centering
  \includegraphics[width=\linewidth]{fig_applications}
  \caption{Start-ups of three conveyors from the literature in physical units, against the
  take-up position (sine profile, $\hat\zeta=0.01$, DIN~22101 strand resistances fitted to
  each $F_U$ and building up with the belt speed, gravity from each carry profile; the take-up
  is moved along the return strand with the case's counterweight). $\sigma_t$: take-up
  position from the head pulley along the return strand, in units of $L$ ($\sigma_t=\xi$ with
  a head drive). Filled triangle: real position; open triangle: alternative of the last row.
  Columns: Lo \citep{lodewijks1996}, \qty{1}{\kilo\metre}, one drive pulley
  ($\egrip=\NumLoGripFactor$), published \qty{30}{\second} start; SM
  \citep{suchorab2025longdistance}, \qty{3}{\kilo\metre} incline, tandem drive
  ($\egrip=16.4$), reference start $3\tper{1}=\qty{\NumSMReferenceStart}{\second}$ (assumed);
  Su \citep{surtees1995runback}, \qty{805}{\metre}, drive on the return strand at
  $\sigma_d=0.197$ ($\egrip=11.5$), published \qty{25}{\second} start. (a--c) Take-up tension
  required by each condition (positive tension along the loop during and after the start;
  grip; running sag below \qty{2}{\percent}), with the case's $\Tt$. (d--f) Shortest sine
  start that the case's $\Tt$ admits, with $3\tper{1}$ and the reference start; grey, no start
  is admissible (the running state already fails). (g) Lo at its real position: the published
  \qty{30}{\second} start (solid) and the shortest admissible one,
  \qty{\NumLoShortest}{\second} (dashed); grip is lost where the entry tension over $\egrip$
  (thin) exceeds the exit tension. (h) SM, real position (solid) and tail (dashed). (i) Su,
  real position (solid) and tight side of the drive (dashed), where the exit tension becomes
  negative. Values in \cref{tab:applications}.}
  \label{fig:applications}
\end{figure}
% Generated by maps/paper_tables.py (phase 5.6) from maps/applications.py; do not edit by hand.
\begin{table*}[t]
\centering
\footnotesize
\setlength{\tabcolsep}{3.5pt}
\begin{threeparttable}
\caption{Start-ups of three conveyors in physical units \citep[Lo, SM, Su:][]{lodewijks1996,suchorab2025longdistance,surtees1995runback}, real take-up position and an alternative (sine profile, $\hat\zeta=0.01$, DIN~22101 strand resistances fitted to $F_U$ and starting with the belt speed, gravity from the carry profile, the case's counterweight in every position). $\sigma_t$: take-up position from the head pulley along the return strand, in units of $L$. Peak drive-entry tension and its ratio to the running value; lowest drive-exit tension during and after the start; carriage travel; take-up tension required by the three conditions (positive tension along the loop, grip $T_\mathrm{entry}\le E\,T_\mathrm{exit}$, running sag below 2\,\%), the one that governs and the running exit tension it implies; shortest sine start that the case's $T_t$ admits.}
\label{tab:applications}
\begin{tabular}{@{}lrrrrrrrrrlrr@{}}
\toprule
 & & & & & \multicolumn{2}{c}{Entry peak} & Exit & & \multicolumn{3}{c}{Required} & Shortest \\
\cmidrule(lr){6-7}\cmidrule(lr){10-12}
Take-up & $\sigma_t$ & $t_1$ & $t_a$ & $t_a/t_1$ & (kN) & /running & min. & Travel & $T_t$ & governs & $T_2$ & start \\
 & & (s) & (s) & & & & (kN) & (m) & (kN) & & (kN) & (s) \\
\midrule
\multicolumn{13}{@{}l}{Lo: 1~km, one drive pulley ($\mathrm{e}^{\mu\theta}=3.0$), $T_t=21.3$~kN; published start 30~s} \\
real, 30~s & 0.001 & 28.5 & 30 & 1.05 & 125 & 2.21 & $21$ & 8.43 & 51.5 & grip & 51.5 & 101 \\
real, shortest start & 0.001 & 28.5 & 101 & 3.56 & 64 & 1.13 & $21$ & 3.20 & 21.3 & grip & 21.3 & 101 \\
tail, 30~s & 0.995 & 24.9 & 30 & 1.20 & 108 & 2.06 & $13$ & 5.15 & 56.1 & grip & 51.8 & -- \\
\midrule
\multicolumn{13}{@{}l}{SM: 3~km incline, tandem drive ($\mathrm{e}^{\mu\theta}=16.4$), $T_t=140$~kN; start $3t_1$ (assumed)} \\
real & 0.100 & 13.6 & 41 & 3.00 & 720 & 1.07 & $156$ & 0.96 & 60.0 & slack & 77.4 & 13 \\
tail & 0.995 & 10.8 & 41 & 3.77 & 726 & 1.07 & $155$ & 0.44 & 21.7 & grip & 46.9 & 10 \\
\midrule
\multicolumn{13}{@{}l}{Su: 805~m, drive at $\sigma_d=0.197$ ($\mathrm{e}^{\mu\theta}=11.5$), $T_t=100$~kN; published start 25~s} \\
real (after the drive) & 0.207 & 6.1 & 25 & 4.11 & 307 & 1.21 & $100$ & 0.25 & 19.8 & sag & 20.4 & 7 \\
tight side & 0.098 & 7.9 & 25 & 3.16 & 95 & 1.01 & $-114$ & 0.00 & 234.0 & grip & 75.2 & -- \\
\bottomrule
\end{tabular}
\begin{tablenotes}[flushleft]
\footnotesize
\item Slack: positive tension along the loop; in SM it is the weight of the return strand held from the take-up near the high end, at rest. Shortest start: --, none (the running state already fails with the case's $T_t$). The shortest start does not account for belt strength or motor torque. Su, tight side: the exit tension is negative already in running.
\end{tablenotes}
\end{threeparttable}
\end{table*}


The formulas of Section~\ref{sec:results} rest on idealised resistances and a horizontal
path. To see how the take-up position acts on real layouts, three conveyors of
\cref{tab:cases} are started in physical units with the full model: a \qty{1}{\kilo\metre}
conveyor with a single drive pulley and published speed-controlled starts
\citep[\S8.6.2]{lodewijks1996}, a \qty{3}{\kilo\metre} incline with a tandem head drive
\citep{suchorab2025longdistance}, and an \qty{805}{\metre} conveyor whose drive sits on the
return strand \citep{surtees1995runback}. Resistances follow DIN~22101 strand by strand, with
the coefficient fitted to the running peripheral force of each case, and build up with the
belt speed; gravity enters through each carry profile; the start is a sine profile with the
base damping $\hat\zeta=0.01$. The take-up is moved along the whole return strand with the
counterweight of the case, so $\Tt$ and $\beta$ stay fixed. At fixed take-up mass the total
tension is $\Tt$ plus a field that does not depend on $\Tt$, so each condition of
Section~\ref{sec:validity} gives a closed bound on $\Tt$ from a single run: positive tension
along the loop during and after the start, drive grip, and a running sag below
\qty{2}{\percent}. The last is the value used in the design calculations of
\citet{surtees1995runback}; DIN~22101 recommends less than \qty{1}{\percent} in steady
operation, mainly for energy efficiency \citep[\S8.2.3]{din22101}. The largest of the three
bounds is the required take-up tension, and the shortest admissible start is the shortest
sine start for which the case's $\Tt$ meets all three, at that duration and beyond. The
second condition, that the take-up follows the belt, never comes close: its margin is at
least \NumAppAccelMarginMin\ in all runs. Resizing the counterweight to keep the running exit
tension, instead of keeping the counterweight, changes the required tension and $\tper{1}$ by
less than \qty{0.3}{\percent}.
% tests: test_applications.py (exact superposition; slow-start limit equal to the running grip
% condition; carriage margin > 10; same-counterweight convention immaterial within 0.3 %)

\paragraph{One drive pulley: the grip sets the start time.} Lodewijks' conveyor has a single
drive pulley with $\egrip=\NumLoGripFactor$ and a take-up tension of
\qty{\NumLoTt}{\kilo\newton}, which he dimensioned for grip with the start-up factor 1.2 of a
drain-type fluid coupling \citep[Table~8.2 and Eqs.~8.7--8.9]{lodewijks1996}. His
\qty{30}{\second} speed-controlled starts last about one fundamental period and fall in the
overshoot of the universal curve: across the start profiles of his Table~8.7 the ratio of
drive-entry to drive-exit tension reaches \NumLoGripRatioMin\ to \NumLoGripRatioMax, against
the limit of \NumLoGripFactor. With the sine profile the belt would need
\qty{\NumLoTtReq}{\kilo\newton} at the take-up, \NumLoSineTtFactor\ times the installed
tension, to stay in grip. His simulations could not show this, since the drive pulleys were
modelled without slip \citep[p.~137]{lodewijks1996}. With his take-up the shortest admissible
sine start is \qty{\NumLoShortest}{\second}, \NumLoShortestOverPeriod\ periods, above the
$2$--$3\,\tper{1}$ of the design rule: with one drive pulley the grip, not the dynamic
amplification, fixes the start time. Moving the take-up away from the drive adds the
resistance of the return run to the required tension, while the required running exit
tension hardly changes (\cref{tab:applications}), and beyond $\xi=\NumLoXiLimit$ the installed
tension no longer holds grip even in steady running (\cref{fig:applications}a,~d). The real
position, next to the drive, is the best one.

\paragraph{Tandem drive on an incline: the statics decide.} The incline of
\citet{suchorab2025longdistance} has a tandem drive ($\egrip=16.4$) and a take-up tension of
\qty{\NumSMTt}{\kilo\newton}. Its take-up sits at $\xi=0.1$, near the high end, where at rest
it must hold the weight of the return strand down the slope: this alone requires
\qty{\NumSMTtReq}{\kilo\newton}, against \qty{\NumSMTtReqTail}{\kilo\newton} with the take-up
at the tail, where grip governs, and up to \qty{\NumSMTtReqMax}{\kilo\newton} next to the
drive. The installed tension covers every position, and the start could be as short as
\qty{\NumSMShortest}{\second}, about \NumSMShortestOverPeriod\ fundamental period. Moving the
take-up to the tail shortens the fundamental period by \qty{\NumSMPeriodChange}{\percent}
(\qty{\NumLoPeriodChange}{\percent} in Lodewijks' conveyor), but changes the peak drive-entry
tension by less than one per cent (\cref{tab:applications}): with a large wave-speed ratio
the carry strand dominates strand~B wherever the take-up is.

\paragraph{Drive on the return strand: the tight side.} Surtees' conveyor has its drive at
$\sigma_d=0.197$ from the head and the take-up right after it, on the slack side. There the
running sag governs and requires \qty{\NumSuTtReq}{\kilo\newton}, close to the lowest value
along the slack side (\qty{\NumSuTtReqBest}{\kilo\newton}); the installed
\qty{100}{\kilo\newton} admits starts from \qty{\NumSuShortest}{\second}, about
\NumSuShortestOverPeriod\ fundamental periods. On the tight side of the drive the take-up
would need \qty{\NumSuTtReqTight}{\kilo\newton}, \NumSuTightFactor\ times more: with the
installed counterweight the exit tension is negative already in running and falls to
\qty{\NumSuExitTight}{\kilo\newton} during the start (\cref{fig:applications}i). The
impractical region of Section~\ref{sec:validity} appears on a real layout.

\paragraph{What the take-up position does on real conveyors.} In these three conveyors the
take-up position acts mainly through the statics: through the resistance of the return run
between the drive and the take-up, and through the weight of the return strand on a slope.
Its dynamic effect, a fundamental period that changes by \NumLoPeriodChange\ to
\qty{\NumSMPeriodChange}{\percent} between the real position and the tail, enters through the
start time, which scales with $\tper{1}$. With one drive pulley the grip fixes the start time;
with a tandem drive or a generous take-up tension the three conditions admit starts of about
one fundamental period, and the limits then come from the belt strength and the motor, which
this analysis does not evaluate. In Lodewijks' and Surtees' conveyors the designers placed
the take-up at, or very near, its best position; in the incline the positions near the tail
would be the best and the one next to the drive the worst, and the installed counterweight
covers both.
